\expandafter\def\csname HistoricalNoteText491\endcsname{\emph{Ueber die Berechnung der Cometenbahnen.} \emph{By} \textsc{Franz Güssmann}. \emph{Wien, bei Joh. Thom. Edlen von Trattnern}, 1803. 4to. 43 pages, with one copperplate. Such a wretched libel, dressed up in the wit of the cowl, would hardly deserve mention at all, were it not worth considering that a man who skulks about solely to discredit deserving men among people who are not properly informed should be shown in his true form and nakedness. This envious and spiteful Jesuit has always written nothing but libels, \emph{à la Valencia}. A small example above shows how conscientiously and perceptively this holy and learned man goes about it. We could give several more if it were worth the trouble.}
\clearpage
\fancyhead[L]{\small GÜSSMANN. COMET ORBITS.}
\hypertarget{native-491-lower}{}\label{native-491-lower}
\begin{center}
\small Monatliche Correspondenz. Vol. XV. pp. 452..460. May 1807.\\[3pt]
\rule{.72\linewidth}{.3pt}
\end{center}
\begin{center}Excerpt from a letter.\end{center}
\begin{flushright}--- --- --- 10 March 1807.\end{flushright}

--- --- --- I have the honour to send you the pamphlet\HistoricalNote{491} by \textsc{Güssmann} that you asked for.%
\SourceLeaf{492}%
\space It hardly deserves a full and detailed reply. Such a reply could be made interesting only by comparing the methods of \textsc{Olbers} and \textsc{Boscovich} and showing the advantages of the former over the latter in ease of use. One can be quite convinced of these advantages without detracting from \textsc{Boscovich}’s fame. Naturally, however, none of that fame benefits his fellow member of the order. The assessment of the latter’s pamphlet, which is in no way suitable for consideration in that investigation, will therefore be better kept separate from it.

To engage with his rude, wretched and altogether irrelevant notions with even a single word would be quite unworthy of you and Dr \textsc{Olbers}. I shall confine myself to picking out the main points on which he might seem to uninformed readers to have said something against Dr \textsc{Olbers}’s method.

\textsc{Güssmann}’s main attack is directed against Dr \textsc{Olbers}’s formula for determining the ratio of the curtate distances in the first and third observations. \textsc{Güssmann} then gives a derivation in \emph{his own manner}, in which he finds a host of absurdities and errors. Dr \textsc{Olbers}, however, will not recognize such a derivation as his own, and anyone who compares the latter’s clear exposition with \textsc{Güssmann}’s peculiar confusion of concepts will find no agreement between them. I do not know whether to regard this pervasive distortion of \textsc{Olbers}’s exposition as \emph{deliberate} or to believe that \textsc{Güssmann} himself failed to understand it. Already in §§ 1 and 2, and later in § 21, the assumption that the motion of the comet is taken to be rectilinear is \emph{falsely attributed} to \textsc{Olbers}. The difference between Dr \textsc{Olbers}’s sole assumption, that the chords of the arcs of the comet’s orbit and the Earth’s orbit traversed between the first and last observations are divided by the middle \emph{radius vectors} in the ratio of the times, and the assumption \textsc{Güssmann} would \emph{foist on him}, that the comet and the Earth move along the chords themselves, is so great that overlooking it, unwittingly or knowingly, takes thoughtlessness or — unparalleled impudence. \textsc{Olbers}’s derivation of the ratio mentioned above ($M$) is as clear as any reader for whom his work is intended could expect. Let $k$ and $i$ be the outer geocentric positions of the comet on the celestial sphere. Let $h$ be the position at which the comet would appear from the Earth if the comet and the Earth were at those points on the chords joining their outer positions where the middle \emph{radius vectors} intersect the chords. Let $h^*$ be the actual geocentric position on the celestial sphere at the middle observation, and let $g$ be the middle position of the Sun. Then $g$, $h$ and $h^*$ lie on a great circle, and so do $k$, $h$ and $i$. The former holds because the hypothetical line of sight evidently lies in the same plane as the Sun and the middle positions of the Earth and the comet. The latter is an easy consequence — one that \textsc{Olbers} rightly left to his readers to work out — of the assumption that the segments into which the chords of the Earth’s orbit and the planet’s orbit are divided by the middle \emph{radius vectors} are proportional to one another. On this basis \textsc{Olbers} devises a simple method for determining the position of $h$ on the celestial sphere. He then uses it to determine the ratio of the curtate distances in the outer observations, $M$. Finally, \textsc{Olbers} shows how one can dispense entirely with the preliminary determination of
\SourceLeaf{493}
$h$ and calculate $M$ directly from the data of the problem. This last step could have been carried out analytically using the formulas already found, but \textsc{Olbers} thought a projection would make the derivation easier to visualize. Now compare this clear exposition with \textsc{Güssmann}’s extraordinary confusion. In § 6 and Fig. II he denotes the outer positions, as I do, by $k$ and $i$, but at first uses $h$ for the \emph{true position}. He claims that \textsc{Olbers} assumes that $h$ lies on the same great circle as $k$ and $i$. This is another shameless falsehood, for \textsc{Olbers} has said no such thing anywhere. Yet immediately afterwards he takes $h$ to be the hypothetical position denoted that way above. This mix-up throws him into complete confusion. He thinks that \textsc{Olbers}’s procedure must always give for $c''$, that is, the longitude of the point $h$, the longitude of the point $h^*$, which is denoted by $a''$. He arrived at this strange judgment only by \emph{arbitrarily} conflating $h$ and $h^*$ and by making the wholly unfounded assumption that the three geocentric positions lie on one great circle. He cannot see at all how the reduction $c''-a''$ could be calculated directly from the observations, because it depends on two elements still unknown, namely the distances of the comet and the Earth from the Sun. This is certainly true. But the explanation is simple. The displacement of the middle position from the great circle through the outer positions also depends on these \emph{very same} elements, and it is precisely this displacement that actually serves to determine $c''-a''$. The fact that the displacement is usually very small makes no difference here, for $c''-a''$ is of the same order. The account \textsc{Güssmann} gives in § 7 of \textsc{Olbers}’s use of the orthographic projection is again entirely wrong. In his Fig. III, the points $b$, $i$ and $r$ have nothing to do with our subject. Only $z$ and $n$ do. The line drawn from $z$ to $n$ relates to the hypothetical position mentioned above, while the line drawn from $b$ to $r$ does not. \textsc{Olbers} denotes the inclination of the projection of the line of sight $z\,n$ by $b''$. It follows that
$\operatorname{tang} b'' = \dfrac{\operatorname{tang} \gamma''}{\sin(A''-c'')}$.
That $\dfrac{\operatorname{tang} \gamma''}{\sin(A''-c'')} = \dfrac{\operatorname{tang} \ObservedBeta''}{\sin(A''-a'')}$ follows at once from the fact that the projections of $z\,n$ and $\mathfrak{S}\,m$ coincide. When \textsc{Güssmann} later suggests in § 9 that \textsc{Olbers}’s formula should really be $\dfrac{t''\sin(a''-a')}{t'\sin(a'''-a'')}$, this is again a consequence of his incomprehensible conflation of the points $h$ and $h^*$. But when he explains the difference between \textsc{Olbers}’s formula and this expression by saying that \textsc{Olbers} includes all three latitudes $\ObservedBeta'$, $\ObservedBeta''$ and $\ObservedBeta'''$, which cannot all fit a straight line at specified points on it as seen from specified points on the Earth’s chord, \textsc{Güssmann} judges quite correctly. Without realizing it, he thereby gives the reason why \textsc{Olbers}’s formula is right, whereas the expression $\dfrac{t''\sin(a''-a')}{t'\sin(a'''-a'')}$, which rests on the assumption of rectilinear motion, is wrong. Indeed, this last formula can be shown to give a useful approximation only when the distances of the Earth and the comet from the Sun are not very unequal — though this is not the place to show it. This happens to be the case in both of \textsc{Olbers}’s examples. It is solely for this reason that the values of $M$ calculated with the expression in question differ from those calculated with \textsc{Olbers}’s formula by no more than \textsc{Güssmann} finds in § 9. The same circumstance accounts for the fact that \textsc{Güssmann} still obtains tolerably accurate results using that expression in his example in § 24. The trivial remarks \textsc{Güssmann} makes in § 11 surely did not escape Dr \textsc{Olbers} either. It is certainly true that $M$ is determined as reliably for small values of $a'''-c''$ and $c''-a'$ as for large ones by $\dfrac{t''\sin(c''-a')}{t'\sin(a'''-c'')}$, provided $c''$ is known with perfect precision. But the limited accuracy of all calculations based on sine tables always leaves some uncertainty in the calculated value of $c''$. \textsc{Olbers} only meant
\SourceLeaf{494}
that this uncertainty, even if it were less than $0''1$, could, in cases where $a'''\!-c''$ and $c''\!-a'$ are very small, distort the value of $M'$. \textsc{Güssmann} is likewise quite right in his second remark. Neglect of the difference between $c''$ and $a''$, which \textsc{Güssmann} denotes by $\delta$, is theoretically inadmissible even when the intervals are as small as can be. But he is wrong to think that $\delta$ makes up a greater part of $a''\!-a'$ as the interval gets smaller. The reverse is true, since $a''\!-a'$ is a quantity of the first order, whereas $\delta$ is of the second order. \textsc{Olbers} knew very well that, even so, neglecting $\delta$ could substantially affect the results even for very small intervals. If he nevertheless allowed this omission in that case, it was because, apart from its absolute magnitude, this effect becomes smaller as the intervals shrink when compared with the inaccuracies arising from unavoidable errors of observation. In cases where the uncertainty of the observations allows only a rough determination of the orbit anyway, it would be quite wrong to make the work more difficult by scrupulously attending to circumstances whose effect is much smaller than that uncertainty and thus to put on a show of precision the circumstances do not permit. As for how far it may be better in such cases to omit the calculation altogether than to derive uncertain results, that is still for discerning astronomers themselves to judge from the particular circumstances.

Here I conclude my remarks. It is enough for me to have discussed the attacks \textsc{Güssmann} made on theoretical grounds. Anyone who cares to compare the brevity and ease of \textsc{Olbers}’s method with the lengthy calculations that \textsc{Güssmann} carries out afterwards, and which nevertheless yield only rough approximations, will also easily form a judgment on this point. The matter speaks for itself. In § 16 \textsc{Güssmann} excuses Mr \textsc{von Pacassi} for the entirely erroneous orbit of the comet of 1779 calculated by \textsc{Boscovich}’s method, on the grounds that in a very \emph{lengthy} calculation it is very easy to make a mistake that then produces an incorrect result --- this is surely telling proof of the advantage Dr \textsc{Olbers}’s method undeniably has in this respect. It is well known that many people, not just one, have readily mastered \textsc{Olbers}’s method and used it to calculate comet orbits reliably and correctly, even people who are neither professional mathematicians nor claim to be particularly practised calculators.
